% acl-header.tex — included via pandoc --include-in-header
%
% Loads ACL conference style on top of pandoc's default LaTeX template.
% Combined with `-V documentclass=article -V classoption=11pt`, this
% reproduces the *ACL paper format (2-column, A4, Times-like font, natbib
% author-year citations).
%
% Used by compile_pdf_pandoc in
% src/academic_pipeline/tools/rendering/latex_compiler.py

% ACL style. `final` (no option) = production; review/preprint variants
% are intentionally omitted because the pipeline ships final-only PDFs.
\usepackage{acl}

% Times-like font and font encoding.
\usepackage[T1]{fontenc}
\usepackage{times}
\usepackage{microtype}

% Tables: pandoc emits longtable for any markdown table > a few rows.
% A separate lua filter (longtable-to-table-star.lua) rewrites those into
% \begin{table*}\begin{tabular}...\end{tabular}\end{table*} so they span
% both columns. The rewritten tabular still uses pandoc's calculated
% column widths (`>{\raggedright\arraybackslash}p{... * \real{0.20}}`)
% which require `array` (for the `>{}` modifier) and `calc` (for
% `\real`). Pandoc's default template loads these only when it thinks
% it needs longtable — and the lua filter strips longtable, so we must
% load them ourselves or the rewritten table* crashes tectonic with
% "Undefined control sequence" on `\real`.
\usepackage{array}
\usepackage{calc}
\usepackage{longtable}
\usepackage{booktabs}

% Figures: graphicx is loaded by pandoc's default template, but we make
% sure caption styling is consistent with ACL conventions (small font,
% italic label).
\usepackage{caption}
\captionsetup[figure]{font=small,labelfont=bf}
\captionsetup[table]{font=small,labelfont=bf}

% Pandoc's tightlist provider is required when --markdown-in-html-blocks
% emits tight unordered lists.
\providecommand{\tightlist}{\setlength{\itemsep}{0pt}\setlength{\parskip}{0pt}}

% Unicode math + punctuation fallbacks: Times (ptmr8t) doesn't ship glyphs
% for these, so remap them to TeX-safe equivalents to avoid "Missing
% character" warnings and \U+FFFD blocks in the output PDF.
\usepackage{newunicodechar}
% Math
\newunicodechar{≥}{$\geq$}
\newunicodechar{≤}{$\leq$}
\newunicodechar{≠}{$\neq$}
\newunicodechar{±}{$\pm$}
\newunicodechar{×}{$\times$}
\newunicodechar{∼}{$\sim$}
\newunicodechar{≈}{$\approx$}
\newunicodechar{→}{$\rightarrow$}
\newunicodechar{←}{$\leftarrow$}
% Punctuation: em-dash, en-dash, smart quotes, ellipsis
\newunicodechar{—}{---}
\newunicodechar{–}{--}
\newunicodechar{‘}{`}
\newunicodechar{’}{'}
\newunicodechar{“}{``}
\newunicodechar{”}{''}
\newunicodechar{…}{\ldots}
% Superscript digits used in p-values like 7.6×10⁻⁶
\newunicodechar{⁰}{$^{0}$}
\newunicodechar{¹}{$^{1}$}
\newunicodechar{²}{$^{2}$}
\newunicodechar{³}{$^{3}$}
\newunicodechar{⁴}{$^{4}$}
\newunicodechar{⁵}{$^{5}$}
\newunicodechar{⁶}{$^{6}$}
\newunicodechar{⁷}{$^{7}$}
\newunicodechar{⁸}{$^{8}$}
\newunicodechar{⁹}{$^{9}$}
\newunicodechar{⁻}{$^{-}$}
